\section{创新性说明}
\subsection{原理创新}
本作品针对传统加密与图像隐写直接结合需要额外组件、神经密码学依赖专门网络训练等问题，
提出了 SSIR——一种由图像密钥控制现有公开 DNN 状态的安全隐蔽信息表示框架。
与已有方案不同，SSIR 不将消息嵌入像素或专用网络，而是将消息表示为目标 DNN 神经元状态。
具体而言，发送方将随机图像与预共享图像密钥融合，输入公开 DNN；
通过分组编码将消息编码为目标比特串，并利用分组神经元响应与目标比特之间的可微损失，
指导对融合图像的对抗优化，直至 DNN 呈现目标状态；随后减去图像密钥得到密文图像。
接收方将密文图像与图像密钥叠加，输入同一 DNN，通过解码神经状态恢复消息。
该原理将密钥控制的消息恢复与隐蔽的公共传输统一起来，密文图像保持自然外观，
非线性响应和概率生成支持输入依赖的密钥控制以及多样的公共表示。
为缓解不响应神经元的影响，本作品设计分组编码，对相邻神经元响应取平均，提高状态诱导可靠性。
此外，本作品识别并抑制两类泄漏：直接视觉泄漏和间接 DNN 状态泄漏，通过基于损失的联合优化同时降低二者。
形式化分析方面，本作品推导了查询有界的 IND-CPA 和 IND-CCA1 优势界，给出了有界查询安全的分布条件和每密钥查询限制。

\subsection{性能创新}
本作品在性能上的创新主要体现在对现有 DNN 基础设施的直接重用和高效生成恢复。
首先，SSIR 无需专门的模型训练，也无需修改所复用 DNN 的参数，原有推理功能得以完整保留，从而显著降低引入专用安全模型的准备与整合开销。
公共 DNN 与编码配置相互独立，为构建独立控制的实现提供了灵活性。其次，本作品通过分组编码和基于损失的优化，在保证可靠性的同时提升生成效率。
实验表明，生成一条 150 位信息的密文图像仅需 300~ms，消息恢复仅需 10~ms，且无需额外模型准备。
SSIR 兼容多种常见 DNN，并在图像失真条件下容忍比特级误差，对多种攻击具有鲁棒性。模块化设计使生成与恢复过程可独立优化，便于在不同部署环境中按需配置。

\subsection{安全与部署创新}
在安全性方面，本作品不仅依赖经验评估，还进行了形式化安全分析。
SSIR 利用高维图像密钥提供灵活的内容表示和大的候选域；
当 KeyGen 提供足够的最小熵时，穷举密钥搜索在计算上不可行。
密钥与传输图像及公共 DNN 配置保持分离，增强了密钥控制的独立性。
形式化分析给出了查询有界的 IND-CPA 和 IND-CCA1 优势界，并规定了有界查询安全所需的分布条件和每密钥查询限制。
这些结论为部署方配置密钥使用与更新策略提供了依据：在允许范围内为每个密钥选择生成和恢复上限，以满足目标优势预算；
计数应在共享同一密钥的会话中累计；在进一步使用可能超过上限之前，各方独立生成并秘密预共享替换密钥，
同时将外部恢复访问限制在已分析的威胁模型内。上述措施符合既定密钥管理指南\supercite{26-barker-2020}，使系统在长期运行中保持可配置的安全边界。
在应用部署上，SSIR 可用于会话密钥的传递或刷新、机密信号通信、
秘密投递通信，以及物联网、应急通信等需要兼顾安全性与隐蔽性的场景，具有较好的实用前景。